В ходе лабораторной работы была решена краевая задача для дифференциального уравнения второго порядка
\[
    y'' - y' = 3
\]
с граничными условиями
\[
    y(0)=6, \qquad y(1)=5e-2.
\]
Для данной задачи было найдено точное аналитическое решение:
\[
    y(x)=1+5e^x-3x.
\]

Для получения численного решения были реализованы и исследованы два метода на равномерной сетке с шагом
\[
    h=0.1.
\]

\begin{itemize}
    \item \textbf{Метод прогонки:} построена конечно-разностная схема для краевой задачи, в результате чего получена трёхдиагональная система линейных алгебраических уравнений. Решение системы найдено методом прогонки с использованием прямого и обратного хода.

    \item \textbf{Метод стрельбы:} применён сеточный вариант метода, при котором решение краевой задачи представляется в виде линейной комбинации двух вспомогательных решений: решения неоднородного уравнения и решения соответствующего однородного уравнения. Для обоих решений использованы рекуррентные формулы, полученные из разностной схемы, после чего неизвестная константа была определена из правого граничного условия.
\end{itemize}

Сравнение численных результатов с точным аналитическим решением показало, что оба метода дают близкие значения и обеспечивают хорошую точность на выбранной сетке. Таким образом, метод прогонки и метод стрельбы были успешно применены для решения поставленной краевой задачи.